% (C) 2026 Ralf Stephan, in collaboration with Claude Code.
% Released under CC0 1.0 Universal (public-domain dedication).
% See https://creativecommons.org/publicdomain/zero/1.0/

\documentclass[11pt]{article}
\usepackage[T1]{fontenc}
\usepackage[utf8]{inputenc}
\usepackage{lmodern}
\usepackage[margin=2.5cm]{geometry}
\usepackage{amsmath,amssymb}
\usepackage{braket}
\usepackage[hidelinks]{hyperref}

% Lean names and paths: verbatim, breakable at . _ /
\DeclareUrlCommand\lean{\urlstyle{tt}}
\def\UrlBreaks{\do\.\do\_\do\/}
% Marks an open problem that blocks checking a main claim of the brief
\newcommand{\blocker}[1]{\textsc{[blocks verification: #1]}~}
% Extra stretch for paragraphs with long, unbreakable Lean names
\emergencystretch=6em

\title{Blocker:\\ Non-linear representations in quantum mechanics}
\author{}
\date{}

\begin{document}
\maketitle

\noindent
This file collects everything about the one item of \lean{open-problems.tex} that blocks
verification, together with the author's replies, a proposed resolution and the changes it needs
in the brief. Sections~\ref{sec:blocker} and~\ref{sec:related} are copied from
\lean{open-problems.tex} and \lean{formalization-log.tex}; the rest is new.

\section{The blocker}\label{sec:blocker}
\begin{enumerate}
\item \blocker{brief issue}\textbf{Theorem~1 as an independence statement.}
  ``REP-VSLIN is independent of QST'' is a statement about logical independence, but
  $\neg$REP-VSLIN is not made precise, and its truth depends on the reading.
  \begin{itemize}
  \item \emph{Literal independence over models.} If a system is only its state space and
    probability, REP-VSLIN follows from QST (take the $H$ that QST provides), so
    $\neg$REP-VSLIN is inconsistent with QST. If a system also carries vector-level data that
    may be arbitrary, any data that is not an inner product space gives a model of
    $\neg$REP-VSLIN, and $m$ is not needed.
  \item \emph{No vector space is fixed, even up to isomorphism.} This is false: by Wigner's
    theorem the Born rule on $P(H)$ fixes $H$ up to a unitary or anti-unitary map (not checked
    in Lean for two different spaces; the Wigner theorem in \lean{ForPhyslib} is for maps of
    one space to itself). For $m(H)$ in particular, $m(H) \cong H$ by construction
    (\lean{twistedIsometry}).
  \item \emph{The linear structure is not fixed by the rays and the Born rule.} This is
    proved (\lean{exists_second_linear_structure}): the carrier of $H$ has a second complex
    inner product space structure, carried over along $m$, with the same scalar
    multiplication, hence the same rays, and the same Born rule, but a different addition. It
    uses eqs.~(3), (4) and (6), and needs $\dim H \ge 2$ (\lean{map_add_of_map_smul}).
  \end{itemize}
  The step of the proof that breaks under the first two readings is ``since the map is
  non-linear, $H \not\equiv m(H)$ in the category of vector spaces'', which is false (Log,
  ``Problems found''). The third reading is proposed as the precise statement. The item stays
  open until the brief adopts it, or another one, and corrects that sentence. (Log,
  ``Differences from the brief''.)
\end{enumerate}

\section{Related items}\label{sec:related}
From \lean{open-problems.tex}, ``From the brief'':
\begin{enumerate}
\setcounter{enumi}{3}
\item \textbf{Any dimension.} The proof needs a non-linear, Born-rule-preserving, ray-preserving
  map on the $\mathcal{H}$ of the system, but constructs $m$ only for a two-state system. The
  same formula $m(\psi) = e^{\imath\theta\,p(z^+|\psi)}\psi$ for a unit vector $\ket{z^+}$
  should work whenever $\dim\mathcal{H} \ge 2$, including infinite dimension; this is not
  formalized ($\mathcal{H} = \mathbb{C}^2$ only). (Brief, Sec.~3, eq.~(2); Log, ``Setting''.)
\end{enumerate}
From \lean{formalization-log.tex}, ``Problems found in the brief'':
\begin{itemize}
\item The closing remark says that $m$ is a diffeomorphism. It is not differentiable at $0$:
  along every real line $m(tv) = t\,m(v)$, so a derivative at $0$ would have to equal $m$
  itself, and $m$ is not additive. The correct statement is that $m$ is a homeomorphism of
  $\mathbb{C}^2$ and a diffeomorphism of $\mathbb{C}^2\setminus\{0\}$. Theorem~1 only needs a
  bijection, so it is not affected.
\item The proof says that ``since the map is non-linear, $H \not\equiv m(H)$ in the category of
  vector spaces''. This is false. As a set $m(H) = H$, since $m$ is a bijection. With the vector
  space structure of $H$, $m(H)$ is $H$. With the structure carried over along $m$, $m$ is by
  construction a linear isometry $H \to m(H)$ (\lean{twistedIsometry}). What is true is that
  the identity of the carrier is not linear from one structure to the other.
\item Eq.~(6) holds for every map that sends each vector into its own ray, since the Born rule
  only sees rays (\lean{born_smul}). Any $v \mapsto \varphi([v])\,v$ with
  $\varphi \colon P(H) \to \mathbb{C}\setminus\{0\}$ would do. What matters for the proof is
  that $m$ commutes with scalars (eq.~(3)) and is not additive (eq.~(4)).
\item QST allows $\dim H \le 1$. There every map that commutes with scalars is additive
  (\lean{map_add_of_map_smul}), so there is no second linear structure. Theorem~1 claims only
  consistency, so the two-state example suffices; ``for every system'' would need
  $\dim H \ge 2$.
\end{itemize}

\section{The author's replies (2026-11-02)}
Paraphrased.
\begin{enumerate}
\item The point is that $H$ and $m(H)$ are still isomorphic, since an isomorphism exists, but
  $m$ is not an isomorphism. --- This is correct when $m(H)$ keeps the operations of $H$: then
  $m$ is not linear. It is equivalent to the third reading above: $m$ is linear from $H$ to
  $H$ exactly when the identity is linear from $H$ to $H$ with the carried-over structure.
\item The formal statement of Theorem~1 is precise but not what matters physically; the
  counterexample is the content. What it means for a physical system to have a mathematical
  characterization cannot be expressed formally.
\item Proposal: the briefs get a section of \emph{mathematical prerequisites}, which states
  what Lean should prove and is fully formalized. The relations between physical conditions
  are argued informally from these prerequisites and are checked by reading, not by Lean. The
  author will rewrite the briefs accordingly.
\end{enumerate}
Under this proposal the blocker is no longer a question of which formal reading of
$\neg$REP-VSLIN is right. It becomes: which mathematical statements go into the prerequisites,
and which premise links them to Theorem~1.

\section{Proposed resolution}
\paragraph{Guess.} The prerequisites state the construction for every complex inner product
space with $\dim H \ge 2$, finite or infinite, and add that it is impossible for
$\dim H = 1$. The informal argument uses the premise:
\begin{quote}
  A mathematical structure of a representation is a physical requirement only if it is fixed
  by the states and their conditional probabilities (the rays and the Born rule).
\end{quote}
Theorem~1 then says, informally: for every system with $\dim H \ge 2$, QST does not fix the
linear structure of $H$, so the linearity in REP-VSLIN is a choice of representation. This is
the third reading of Section~\ref{sec:blocker}, with the premise made explicit.

\paragraph{Why.}
\begin{itemize}
\item \emph{The premise is the brief's own.} The abstract says that $m$ ``preserves the Born
  rule, and therefore all the physics''. This is operational equivalence: two descriptions
  that give the same probabilities for every measurement are the same physics. It is also the
  premise of Wigner's theorem.
\item \emph{$m$ changes only unobservable phases.} $m$ multiplies each vector by a phase that
  depends only on its ray. The overall phase of a vector cannot be measured, so on states $m$
  is the identity: a relabelling, not a new theory.
\item \emph{Wigner's theorem leaves exactly this gap.} A bijection of $P(H)$ that preserves the
  Born rule comes from a unitary or anti-unitary map, but this concerns rays. Nothing forces the
  map on vectors to be linear, and $m$ induces the identity on rays while being non-linear on
  vectors. The gap exists in every dimension $\ge 2$, so the prerequisite should be general.
\item \emph{The conclusion is meant for every system.} ``Linearity is not part of the
  physics'' is a claim about quantum systems in general, not about qubits. One example
  suffices for logical consistency, but not for the intended claim.
\item \emph{Dimension $1$ is not a physical system.} It has one state, its only observables are
  multiples of the identity, and every probability is $1$. Excluding it costs nothing, and it
  is exactly where the construction fails (\lean{map_add_of_map_smul}).
\item \emph{Superposition is not lost.} Interference lies in the probabilities between rays,
  which $m$ preserves. Superposition in $m(H)$ is the carried-over addition, and every
  linear operator $U$ on $H$ becomes the operator $mUm^{-1}$, linear for it. Unlike non-linear
  dynamics in the style of Weinberg, which allows signalling, this predicts nothing new.
\end{itemize}

\paragraph{Confidence.}
\begin{itemize}
\item High: the premise above is what the author means; the prerequisites should hold for
  $\dim H \ge 2$ in general, with $\dim H = 1$ excluded.
\item Medium: the scope is a single system. For composite systems, applying $m$ to each
  factor while keeping the usual tensor product need not be consistent; the tensor product
  would have to be carried over along $m$ as well. The brief does not discuss composite
  systems, so it should either say ``a single system'' or treat them.
\end{itemize}

\section{Mathematical prerequisites to add to the brief}
Each item is a purely mathematical statement for Lean. Let $H$ be a complex inner product space,
$z \in H$ a unit vector and $\theta \in \mathbb R$. Write
$p(z\mid v) = |\braket{z|v}|^2/\|v\|^2$ for $v \ne 0$, and define $m_\theta(0) = 0$ and
$m_\theta(v) = e^{\imath\theta\,p(z\mid v)}\,v$ for $v \neq 0$.
\begin{enumerate}
\item[P1.] \emph{Properties of $m_\theta$.} $m_\theta$ is a homeomorphism of $H$ with inverse
  $m_{-\theta}$, and a diffeomorphism of $H\setminus\{0\}$. It commutes with every complex
  scalar, sends each vector into its own ray with a factor of modulus $1$, and preserves norms
  and the Born rule. It is not differentiable at $0$ when $\dim H \ge 2$ and
  $\theta \in (0,2\pi)$.\par{\raggedright
  Lean, for $H = \mathbb C^2$ only: \lean{nlMap_smul}, \lean{nlMap_mem_ray},
  \lean{norm_nlMap}, \lean{born_nlMap}, \lean{nlHomeomorph}, \lean{contDiffAt_nlMap},
  \lean{not_differentiableAt_nlMap_zero}.\par}
\item[P2.] \emph{Non-linearity.} If $\dim H \ge 2$ and $\theta \in (0, 2\pi)$, then $m_\theta$
  is not additive.\\
  Lean, for $H = \mathbb C^2$ only: \lean{nlMap_not_add}, \lean{not_isLinearMap_nlMap}.
\item[P3.] \emph{A second linear structure.} If $\dim H \ge 2$ and $\theta \in (0, 2\pi)$, the
  set $H$ carries a second complex inner product space structure, carried over along
  $m_\theta$, with the same scalar multiplication (hence the same rays) and the same Born rule,
  but a different addition. $m_\theta$ is a linear isometry from $H$ onto it.\\
  Lean, for $H = \mathbb C^2$ only: \lean{exists_second_linear_structure},
  \lean{twistedIsometry}.
\item[P4.] \emph{Dimension one.} If $\dim H = 1$, every map $H \to H$ that commutes with complex
  scalars is additive, so no second linear structure with the same scalar multiplication
  exists.\\
  Lean, for every $1$-dimensional space: \lean{map_add_of_map_smul}.
\end{enumerate}
Lean work needed: P1--P3 for a general $H$ with $\dim H \ge 2$ (item~4 of
Section~\ref{sec:related}). The formula is the same; the two-state proofs should generalize
with $z^+$ replaced by $z$ and $z^-$ by a unit vector orthogonal to $z$.

\section{Other changes to the brief}
\begin{enumerate}
\item \emph{State the premise.} Add the premise of the previous section, for instance after
  Condition REP-VSLIN, as the link between the prerequisites and Theorem~1.
\item \emph{Theorem~1.} Restate it as the informal consequence of P3 and the premise: for a
  system with $\dim H \ge 2$, the states and the Born rule do not fix the linear structure of
  $H$. If ``independent'' stays, say what $\neg$REP-VSLIN means in these terms.
\item \emph{The sentence $H \not\equiv m(H)$.} Replace ``Since the map is non-linear,
  $H \not\equiv m(H)$ in the category of vector spaces'' by, for example: ``With the operations
  carried over along $m$, the set $H$ is again an inner product space, isomorphic to $H$ via
  $m$; but $m$ is not linear for the operations of $H$, so the same states and Born rule are
  represented by two different linear structures.''
\item \emph{Dimension.} Replace ``Consider a two-state system'' by a general $H$ with
  $\dim H \ge 2$ and a unit vector $\ket{z^+}$, keeping the two-state case and Fig.~1 as the
  illustration. Note that for $\dim H = 1$ no such map exists (P4).
\item \emph{Diffeomorphism.} In the Remark, ``$m$ is \ldots\ a diffeomorphism'' should read:
  a homeomorphism of $H$ and a diffeomorphism of $H \setminus \{0\}$ (P1).
\item \emph{The Remark on transport.} It can cite P3: since $m$ is a linear isometry onto the
  carried-over structure, superpositions, observables and time evolution transfer along $m$.
  This settles item~1 of \lean{open-problems.tex} for everything defined from the linear and
  inner product structure.
\item \emph{Scope.} Say ``a single system'', or treat composite systems (see Confidence).
\item \emph{Optional.} Note that eq.~(6) holds for any map that rescales within rays; what the
  proof needs is eq.~(3) and eq.~(4).
\end{enumerate}

\paragraph{When the blocker closes.} Once the brief has the prerequisites P1--P4 and the
premise, Lean verifies P1--P4, and Theorem~1 is recorded in the log as argued informally from
them and checked by reading.

\end{document}
